% !TEX program = xelatex
\documentclass[aspectratio=169,11pt]{beamer}
\usepackage{fontspec}
\usepackage[brazilian]{babel}
\usepackage{booktabs,graphicx,tikz}
\setmainfont{Noto Sans}
\setsansfont{Noto Sans}
\renewcommand{\familydefault}{\sfdefault}
\graphicspath{{../figuras/}}

% --- paleta: tinta, cinza e um único acento (o mesmo azul das figuras)
\definecolor{ink}{HTML}{1B1F24}
\definecolor{muted}{HTML}{6B7280}
\definecolor{accent}{HTML}{1F5FA8}
\definecolor{warn}{HTML}{C0392B}
\setbeamercolor{normal text}{fg=ink,bg=white}
\setbeamercolor{structure}{fg=accent}
\setbeamercolor{itemize item}{fg=accent}
\setbeamertemplate{itemize item}{\raisebox{.15ex}{\scriptsize$\bullet$}}
\setbeamertemplate{navigation symbols}{}
\setbeamersize{text margin left=10mm,text margin right=10mm}
\setbeamerfont{frametitle}{size=\large,series=\bfseries}
\setbeamertemplate{frametitle}{\vskip5mm\raggedright\color{ink}\usebeamerfont{frametitle}\insertframetitle\par\vskip2mm}
\setbeamertemplate{footline}{\hfill{\scriptsize\color{muted}\insertframenumber\hspace{10mm}}\vskip4mm}

% número grande + rótulo
\newcommand{\stat}[3][accent]{\begin{minipage}[t]{.3\textwidth}\raggedright{\fontsize{26}{28}\selectfont\bfseries\color{#1}#2}\\[1mm]{\small\color{muted}#3}\end{minipage}}
\newcommand{\fig}[2][.7]{\includegraphics[width=\linewidth,height=#1\textheight,keepaspectratio]{#2}}
\newcommand{\nota}[1]{{\small\color{muted}#1}}

\begin{document}

% ---------- 1. capa
\begin{frame}[plain]
\vspace*{18mm}
{\fontsize{22}{26}\selectfont\bfseries Caxias Heavy Industries:\\[1mm] como sair do vermelho\\ em 12 meses\par}
\vspace{4mm}\color{accent}\rule{22mm}{1.2pt}\color{ink}\par\vspace{4mm}
{\large Plano de reestruturação com sistema multiagentes}\par
\vspace{14mm}
{\small\color{muted}
Iuri, João Guedes, Bruno Muniz, Davi, Diogo e Novak\\ IA Generativa, UFSM \;·\; Novembro de 2026}
\end{frame}

% ---------- 2. equipe
\begin{frame}{Quem fez o quê}
\vspace{2mm}
\begin{tabular}{@{}p{.38\textwidth}p{.45\textwidth}@{}}
\toprule
\textbf{Papel} & \textbf{Integrante}\\
\midrule
Consultor, contato único com o cliente & Iuri
\\[2mm]
Scrum Master & João Guedes
\\[2mm]
Engenharia de dados & Bruno Muniz
\\[2mm]
Desenvolvimento dos agentes & Davi
\\[2mm]
Análise e validação & Diogo e Novak
\\
\bottomrule
\end{tabular}
\end{frame}

% ---------- 3. diagnóstico geral
\begin{frame}{A empresa tem pedidos, mas não tem folga}
\vspace{3mm}
\stat{273,9}{R\$ milhões: carteira de 12 meses (320 ônibus)}\hfill
\stat{4,36}{R\$ milhões por mês: lucro bruto das entregas}\hfill
\stat[warn]{3,77}{R\$ milhões por mês: dividendos do plano (corrente + atrasados em 24 meses)}
\vspace{5mm}

{\large O lucro bruto cobre os dividendos \textbf{1,16 vez}.} Com caixa zero, qualquer atraso vira rombo.

\vspace{4mm}\nota{Dividendos de 12 meses: R\$ 33,3 mi, dos quais R\$ 22,5 mi exigíveis e R\$ 10,8 mi de ações ON suspensas.}
\end{frame}

% ---------- 4. margem
\begin{frame}{Todos os modelos dão lucro; o elétrico dá o maior por ônibus}
\begin{columns}[c]
\column{.62\textwidth}\fig[.68]{02_margem_modelos.png}
\column{.34\textwidth}
Margem de 25,5\% (elétrico) a 32,3\% (rodoviário) no preço praticado.\\[3mm]
Lucro por ônibus: de R\$ 109 mil (micro) a R\$ 496 mil (elétrico).
\end{columns}
\end{frame}

% ---------- 5. pendentes
\begin{frame}{27 ônibus pendentes seguram R\$ 28 mi}
\begin{columns}[c]
\column{.6\textwidth}\fig[.68]{03_carteira_pedidos.png}
\column{.36\textwidth}
{\fontsize{24}{26}\selectfont\bfseries\color{warn}27}\\[-1mm]{\small\color{muted}pedidos pendentes ou atrasados}\\[4mm]
Prefeitura só paga na entrega.\\[2mm]
Multa de R\$ 2,81 mi mais R\$ 28 mil por dia.
\end{columns}
\end{frame}

% ---------- 6. capacidade
\begin{frame}{No papel a linha dá conta; o gargalo está fora dela}
\begin{columns}[c]
\column{.64\textwidth}\fig[.68]{04_capacidade_linha.png}
\column{.32\textwidth}
89 ônibus/mês com 3 turnos; 59 com 2. O pico de pedidos é 46.\\[3mm]
Os atrasos vêm de fornecedores e paradas de máquina, que a planilha não mede.
\end{columns}
\end{frame}

% ---------- 7. refugo
\begin{frame}{Refugo importa pouco no caixa}
\begin{columns}[c]
\column{.62\textwidth}\fig[.68]{05_refugo_etapas.png}
\column{.34\textwidth}
Etapas 20, 50 e 60 passam de 5\%.\\[3mm]
Só 71\% dos ônibus passam sem retrabalho, mas o retrabalho é 4\% do tempo de ciclo.
\end{columns}
\end{frame}

% ---------- 8. fornecedores
\begin{frame}{O inversor chinês prende R\$ 3,9 mi no caminho}
\begin{columns}[c]
\column{.62\textwidth}\fig[.68]{06_leadtime_pecas.png}
\column{.34\textwidth}
Inversor: 60 dias, sem produtor nacional. Chicote Yazaki: não substituível.\\[3mm]
Dólar +10\% custa R\$ 245 mil por mês.
\end{columns}
\end{frame}

% ---------- 9. arquitetura
\begin{frame}{Sete agentes, com um revisor que pode vetar}
\begin{center}\fig[.78]{01_arquitetura.png}\end{center}
\end{frame}

% ---------- 10. revisor
\begin{frame}{O Revisor recalcula e aplica as regras do cliente}
\begin{columns}[T]
\column{.48\textwidth}
\textbf{Por que multiagentes}
\begin{itemize}
  \item Quatro domínios independentes
  \item Dados externos do Banco Central
  \item Contas feitas em Python, não pelo LLM
\end{itemize}
\column{.48\textwidth}
\textbf{O que o Revisor impõe}
\begin{itemize}
  \item Horizonte de 12 meses
  \item Rescisão CLT; aviso de 30 dias a fornecedor
  \item Chicote e inversor ficam
  \item Caixa zero
\end{itemize}
\end{columns}
\vspace{3mm}
\hrule\vspace{3mm}
\textbf{\color{warn}Exemplo de veto:} reduzir o elétrico perde R\$ 12,2 mi de lucro para liberar R\$ 2,0 mi. Reprovado.
\end{frame}

% ---------- 11. plano
\begin{frame}{11 propostas: 1 reprovada, 2 condicionais}
\footnotesize
\begin{tabular}{@{}l@{\hspace{3mm}}l@{\hspace{3mm}}r@{\hspace{3mm}}l@{}}
\toprule
 & \textbf{Proposta} & \textbf{Efeito} & \textbf{Parecer}\\
\midrule
FIN-1 & Dividendos: corrente em dia, atrasados em 24 meses & $-$R\$ 0,99 mi/mês & aprovada\\
FIN-2 & Entregar os 27 pendentes & +R\$ 7,07 mi (único) & aprovada\\
FIN-3 & Reduzir a produção do elétrico & & \color{warn}reprovada\\
FIN-4 & Diretoria $-$20\% (parecer jurídico) & +R\$ 77 mil/mês & condicional\\
OPE-1 & 2 turnos e ajuste do quadro & +R\$ 58 mil/mês & condicional\\
OPE-2 & Manutenção nas etapas 20, 50 e 60 & +R\$ 2 mil/mês & aprovada\\
SUP-1 & Comprar peças longas no pedido firme & +R\$ 421 mil (único) & aprovada\\
SUP-2/3 & Cotar alternativas; elétrico sob pedido & não somado & aprovada\\
MER-1/2 & Mix de maior lucro; double-deck na fase 2 & não somado & aprovada\\
\bottomrule
\end{tabular}
\end{frame}

% ---------- 12. resultado
\begin{frame}{Com o plano, a cobertura sobe e o saldo fica positivo}
\begin{columns}[c]
\column{.64\textwidth}\fig[.68]{07_caixa_dividendos.png}
\column{.32\textwidth}
{\fontsize{18}{20}\selectfont\bfseries\color{accent}1,16x $\to$ 1,19x}\\[-1mm]{\small\color{muted}cobertura dos dividendos}\\[4mm]
Saldo mínimo de R\$ 1,0 mi no mês 1.\\[1mm]
Selic sobre os atrasados: R\$ 1,29 mi.\\[2mm]
\nota{Saldo é proxy: lucro bruto menos dividendos.}
\end{columns}
\end{frame}

% ---------- 13. limites
\begin{frame}{O que ainda precisamos confirmar}
\begin{columns}[T]
\column{.48\textwidth}
\textbf{Limitações}
\begin{itemize}
  \item O saldo não inclui outras despesas nem dívidas
  \item Capacidade teórica, sem paradas de máquina
  \item Operadores por estação, rescisão e multa são hipóteses
  \item O modo com LLM ainda não foi executado
\end{itemize}
\column{.48\textwidth}
\textbf{Próximos passos}
\begin{itemize}
  \item Juros e multa dos dividendos atrasados
  \item Dias de atraso e custo do refugo
  \item Cotar peças alternativas
  \item Fase 2: double-deck e P\&D, depois do verde
\end{itemize}
\end{columns}
\end{frame}

\end{document}
